% ======================= 第 一 章 =======================
\chapter{两座城}
\chapterkana{二つの都}
\epigraphbox{“月出皎兮，佼人僚兮。”}{《诗经·陈风·月出》}

\kaishou{月之都不是一座宫殿，是一座城，而且有两半：一半在背面，是本体；一半在正面，是伸出去的一角。本章要说明这两半各自的样貌，以及它们之间隔着什么。}

\noindent
「月之都」这三个字，通行的理解是一座宫殿：琉璃为瓦，白玉为阶，住着一位不能回头的公主。

这个理解有一半错。宫殿是有的，只是不止一座；住的人也不止一位。真正的问题在于「一座」这个词：它让人以为那地方是一个点，实际上它是一个面，正反两面各有一处，中间隔着一道从外面看不见的界线。

\section{背面}

月球被潮汐锁定，永远以同一面朝向地球。这不是一个新知识，但它有一个后果常被忽略：背面不是「看不见的那一面」，是一个从来没有被地面观察过的半球。

一九五九年十月七日，苏联的「月球三号」在飞越月球时第一次拍到它，二十九帧，最清楚的一帧一个像素抵得上一公里。此后的六十年里，人类又拍了十几次，分辨率从一公里提到半米。二〇一九年一月，「嫦娥四号」在背面的冯·卡门坑落下去，就地测了三个月。

这些影像给出一个结论，与本书的题目直接相关：背面的地形可以归成两类，一类是高地的坑洼，占绝大部分；另一类很少，不到百分之三，平坦，反照率低，边缘等宽。

等宽到什么程度，可以说一个数：沿边界取一百二十个测点，曲率半径的标准差是百分之三；同一批数据里，高地边缘的对应值是百分之十九。

自然的边界做不出这种整齐。撞击坑随地形起伏，熔岩流随黏度变化，风成堆积随粒径分选。三片低地的整齐程度，比测量本身的噪声还低。

还有一条更硬的证据：重力场测量显示，那几片低地底下没有壳厚异常。它不是盆地，也不是熔岩填出来的凹地。形状是平的，地质上没有理由。

\begin{figure}[htbp]
\centering
\begin{tikzpicture}[x=0.88mm,y=1mm]
  \useasboundingbox (-6,-13) rectangle (156,26);
  \draw[ink,line width=0.4pt] (0,0) -- (148,0);
  \draw[ink,line width=0.4pt] (0,0) -- (0,22);
  \node[anchor=north,font=\scriptsize] at (0,-1.5) {0};
  \node[anchor=north,font=\scriptsize] at (74,-1.5) {60};
  \node[anchor=north,font=\scriptsize] at (148,-1.5) {120 km};
  \node[anchor=east,font=\scriptsize] at (-2,11) {反照率};
  \draw[rulegrey,line width=0.4pt]
    plot[domain=0:36,samples=70] ({\x},{12+1.6*sin(\x*57)});
  \draw[ink,line width=0.6pt] (36,13.1) -- (110,13.0);
  \draw[ink,line width=0.6pt] (110,13.0) -- (146,12.0);
  \draw[ink,line width=0.9pt,dashed] (36,2) -- (36,20);
  \node[anchor=south,font=\scriptsize] at (36,20) {界线};
  \node[anchor=north west,font=\scriptsize,align=left] at (40,10)
    {平滑区：标准差 0.0016};
  \node[anchor=north west,font=\scriptsize,align=left] at (2,8)
    {高地：标准差 0.011};
\end{tikzpicture}
\caption{背面某低地的东西向反照率剖面（重绘）}
\figsource{据克莱门汀一九九四年反照率数据重绘，纵轴已归一化。原始采样间隔 480 m，此处按 1 km 重采样。}
\label{fig:albedo}
\end{figure}

图 \ref{fig:albedo} 横穿一片低地。左侧是高地的起伏，右侧拉平，中间那段虚线就是界线。要紧的细节在界线处：曲线没有过渡段，从起伏到平坦隔了不到两个采样点。

三条证据合起来，只能排除自然成因，不能证明人工成因。要说成人工，还缺一样东西，就是那圈边界本身。影像上看不出它的厚度，测高看不出它的高低，雷达只探到两米。

\section{结界}

把这几片低地连起来看，它们占的是一个半球的内侧。本书要提出的读法很朴素：那不是几片平地，是一个内壁。

内壁的说法有一个直接后果：壁内壁外不是一个环境。这正好解释了背面影像里一处长期无法解释的异常：反照率的整体偏低。

低到什么程度，可以给个数：背面高地的反照率在 0.11 到 0.13 之间，与正面高地一致；而那几片低地在 0.07 上下，比月海还暗。暗的原因通常是有机质或细粒物质，月面上两者都不该有。

本书的解释是：那层暗色不是地表，是壁内的沉积。壁内的空气把尘埃托起来，慢慢沉在低处，两万年下来积成一层膜。

壁内是什么样子，只有一处间接材料。一七六三年，一艘法国商船在好望角外海捞起一只铅盒，内藏铜牌一枚，径四十一毫米，两面刻有同一幅图：一个圆，圆内三条弧线，弧线之间写着一种无人能识的小字。铜牌一九六三年入 `ФВ` 卷宗，编号 `ФВ-1963-44`。

那三条弧线，与背面影像里三片低地的排布一致。

\begin{archive}[机密／СЕКРЕТНО]{档案 № ФВ-1963-44}
铜牌一枚，直径 41 mm，厚 2.7 mm，铜锌比为 82 比 18，与十八世纪欧洲黄铜不合。两面同图，边缘有等分刻度 108 格。

盒内另附纸片一张，已脆，仅存半行，墨书，字体为十八世纪法文手写体，可辨者为「……nous n'avons pas pu entrer」，意为「我们没能进去」。

纸片背面另有一行铅笔字，字迹与正面不同，写作：「他们也一样。」
\end{archive}

档案里那半行法文，是本书能找到的最早一句关于那道界线的记录。它不说明界线是什么，只说明一七六三年有人到过那里，并且没有进去。

\section{桃园与营造}

壁内有一座城，这一点是从另一类材料推的：壁内大气的光谱。

二〇〇九年以后，几台探测器在背面低地上空测到了极微弱的稀薄气体，主要成分是氮与氧，浓度约为地球海平面的十万分之一。这个量太小，不足以形成大气层，却足以说明一件事：那些气体在往外漏。

漏，就有一个来源。本书认为来源是城本身：一座维持了两万年的人工环境，其废气的排出速率与测得的漏失量在同一个数量级。

城里有三样东西已经被间接证实：一处人工林地、一套营造法式、一条通行规则。

人工林地，他们自己叫桃园。它不是农地，不产粮，用途是供氧与休憩。这一条的根据是影像上的成片纹理：在最大的那片低地北缘，有一块约十四平方公里的区域，纹理呈规则的条带状，间距 9.4 米，与任何一种自然植被都不符。条带的方向与低地的长轴一致，说明它是按地形的走向排的。

营造法式是另一回事。城里建筑的形制为中国风格，这一点在影像上无法直接看出，是从两处旁证推的：一是撤离之后落地的那批人所建的最早一批聚落，其夯土基址的模数与中原仰韶晚期的做法一致；二是他们带下来的营造手册在后世被引用过，引用的条目里有「柱高一丈二，础方三尺」这类具体的模数。

通行规则只有一条，且很硬：那道界线只限通行，不限信号。

\section{两半}

正面那一处与背面这一座，是同一个东西的两半。

正面的一处在静海西缘，编号 MARE-TRANQ-07，中文材料里叫广寒宫。它的规模比背面那座城小得多：外有一道屏蔽环，直径约三点五公里；环内是一道垣，周长九里，约合一点二公里；垣内有六道同心环，环与环之间有门，门不同时开启。

背面那座城的规模，从低地的面积反推：三片低地合计约六万平方公里，是正面那处的四千倍。

两半之间隔着那道界。界的性质，本书只能给三条：与月船同龄；只限通行；前二二一年之后，随全城一并失效。

前二二一年这个年份，是本书后半部分的起点，也是中国史的转折点。两件事发生在同一年，不是巧合，第七章会说明。

\begin{table}[htbp]
\centering
\caption{两半的规模}
\label{tab:two}
\small
\begin{tabular}{@{}lll@{}}
\toprule
项 & 正面（静海七号） & 背面（本体）\\
\midrule
位置 & 静海西缘，正面 & 静海正对面，背面\\
外缘直径 & 约 3.5 km（屏蔽环） & 约 6 万 km$^2$（三片低地合计）\\
垣 & 周长九里，高十二丈 & 未见\\
环制 & 六道同心环 & 未见\\
空气 & 无（在地面环境里） & 有，浓度约地球十万分之一\\
通行 & 仅经垣门 & 仅经那道界线\\
\bottomrule
\end{tabular}
\tabsource{正面一栏据《广寒宫残简》辑校本与《战争的雷霆》第一章；背面一栏据本书第二章所论。两栏均为推算值。}
\end{table}

表 \ref{tab:two} 里最后一行，是本书最没有把握的一处。正面那处的通行规则有简牍为证，背面那处的只有一句档案里的法文：「我们没能进去」。

笔者把两半并成一张表，是要说明一件事：这座城的形制是自上而下统一的，正面那处是它伸出来的一个角。可这个判断有一半建立在一个负数上：我们进不去，所以我们不知道里面。

这个负数会一直带到全书结束。
